% GENERATED FILE. Edit canonical lessons, then run npm run generate:books.
% canonical-source-hash: fnv1a64:de5ac37589287e01
% canonical-lessons: TA-C462-avamanam, TA-C462-cankatam, TA-C462-kuccam, TA-C462-aruvaruppu, TA-C462-mukam-culi

\chapter{Shame and Disgust}
\label{ch:ta-shame-and-disgust}
\hlchaptermodality{462}

\begin{chapteropening}
\textbf{By the end of this chapter:} \emph{I can talk about humiliation, embarrassment, shyness and disgust.}

Complete the last lesson of chapter 462: I can talk about humiliation, embarrassment, shyness and disgust.
\end{chapteropening}

\section[\texorpdfstring{\ta{அவமானம்} (avamāṉam)}{அவமானம் (avamāṉam)}]{\ta{அவமானம்} (avamāṉam) — humiliation; disgrace}

\label{lesson:TA-C462-avamanam}



\begin{quote}
Before the new one: say the Tamil for keenness, then the Tamil for eagerness.
\end{quote}

\subsection*{You'll want to know: \ta{அவமானம்}}
\textbf{\ta{அவமானம்}} — \emph{avamāṉam} — \textquotedblleft{}humiliation; disgrace\textquotedblright{}.

\begin{quote}
\ta{எல்லோர்} \ta{முன்னாலும்} \ta{திட்டியதால்} \ta{அவனுக்கு} \ta{அவமானமாக} \ta{இருந்தது}.
\end{quote}

\begin{quote}
\emph{ellōr muṉṉālum tiṭṭiyatāl avaṉukku avamāṉamāka iruntatu.}
\end{quote}

\begin{quote}
He felt humiliated because he was scolded in front of everyone.
\end{quote}

\ta{அவமானப்படுத்து}, to humiliate someone.

\begin{cousinweb}
\textbf{\ta{அவமானம்}} (\emph{avamāṉam}) is borrowed from Sanskrit \emph{avamāna}, \textquotedblleft{}disrespect, contempt\textquotedblright{} (Monier-Williams).
\end{cousinweb}

\begin{grammarlens}[title={What you've built}]
One more word for describing what you lived and felt.
\end{grammarlens}

\subsection*{Your turn}
\begin{itemize}
\raggedright
  \item \emph{Say it:} \emph{avamāṉam}
  \item \emph{Say it:} \emph{avamāṉam}, once more
  \item \emph{Say it:} \emph{ellōr muṉṉālum tiṭṭiyatāl avaṉukku avamāṉamāka iruntatu.} — He felt humiliated because he was scolded in front of everyone.
  \item \emph{Recall:} say the Tamil for keenness, then the Tamil for eagerness, then say \emph{avamāṉam} again
\end{itemize}

\begin{tcolorbox}[breakable,colback=teal!4,colframe=teal!35!black,title={Before you move on}]
What does \ta{அவமானம்} mean? (\textquotedblleft{}Humiliation; disgrace\textquotedblright{}.) Say it once more.
\end{tcolorbox}
\section[\texorpdfstring{\ta{சங்கடம்} (caṅkaṭam)}{சங்கடம் (caṅkaṭam)}]{\ta{சங்கடம்} (caṅkaṭam) — awkwardness; embarrassment}

\label{lesson:TA-C462-cankatam}



\begin{quote}
Before the new one: say the Tamil for eagerness, then the Tamil for humiliation.
\end{quote}

\subsection*{You'll want to know: \ta{சங்கடம்}}
\textbf{\ta{சங்கடம்}} — \emph{caṅkaṭam} — \textquotedblleft{}awkwardness; embarrassment\textquotedblright{}.

\begin{quote}
\ta{அவருடைய} \ta{கேள்வி} \ta{எனக்குச்} \ta{சங்கடமாக} \ta{இருந்தது}.
\end{quote}

\begin{quote}
\emph{avaruṭaiya kēḷvi eṉakkuc caṅkaṭamāka iruntatu.}
\end{quote}

\begin{quote}
His question made me uncomfortable.
\end{quote}

\ta{சங்கடமான} \ta{நிலை}, an awkward situation.

\begin{cousinweb}
\textbf{\ta{சங்கடம்}} (\emph{caṅkaṭam}) is borrowed from Sanskrit \emph{saṁkaṭa}, \textquotedblleft{}narrow; a strait, a difficulty\textquotedblright{} (CDIAL 12817).
\end{cousinweb}

\begin{grammarlens}[title={What you've built}]
One more word for describing what you lived and felt.
\end{grammarlens}

\subsection*{Your turn}
\begin{itemize}
\raggedright
  \item \emph{Say it:} \emph{caṅkaṭam}
  \item \emph{Say it:} \emph{caṅkaṭam}, once more
  \item \emph{Say it:} \emph{avaruṭaiya kēḷvi eṉakkuc caṅkaṭamāka iruntatu.} — His question made me uncomfortable.
  \item \emph{Recall:} say the Tamil for eagerness, then the Tamil for humiliation, then say \emph{caṅkaṭam} again
\end{itemize}

\begin{tcolorbox}[breakable,colback=teal!4,colframe=teal!35!black,title={Before you move on}]
What does \ta{சங்கடம்} mean? (\textquotedblleft{}Awkwardness; embarrassment\textquotedblright{}.) Say it once more.
\end{tcolorbox}
\section[\texorpdfstring{\ta{கூச்சம்} (kūccam)}{கூச்சம் (kūccam)}]{\ta{கூச்சம்} (kūccam) — bashfulness; ticklishness}

\label{lesson:TA-C462-kuccam}



\begin{quote}
Before the new one: say the Tamil for humiliation, then the Tamil for awkwardness.
\end{quote}

\subsection*{You'll want to know: \ta{கூச்சம்}}
\textbf{\ta{கூச்சம்}} — \emph{kūccam} — \textquotedblleft{}bashfulness; ticklishness\textquotedblright{}.

\begin{quote}
\ta{புதியவர்களிடம்} \ta{பேச} \ta{என்} \ta{மகனுக்குக்} \ta{கூச்சம்}.
\end{quote}

\begin{quote}
\emph{putiyavarkaḷiṭam pēca eṉ makaṉukkuk kūccam.}
\end{quote}

\begin{quote}
My son is shy about talking to new people.
\end{quote}

Also being ticklish: \ta{கூச்சமாக} \ta{இருக்கிறது}, it tickles.

\begin{cousinweb}
\textbf{\ta{கூச்சம்}} (\emph{kūccam}) is native Dravidian, from \ta{கூசு} \emph{kūcu}, \textquotedblleft{}to be shy, to be ticklish, to shrink back\textquotedblright{} (DEDR 1876).
\end{cousinweb}

\begin{grammarlens}[title={What you've built}]
One more word for describing what you lived and felt.
\end{grammarlens}

\subsection*{Your turn}
\begin{itemize}
\raggedright
  \item \emph{Say it:} \emph{kūccam}
  \item \emph{Say it:} \emph{kūccam}, once more
  \item \emph{Say it:} \emph{putiyavarkaḷiṭam pēca eṉ makaṉukkuk kūccam.} — My son is shy about talking to new people.
  \item \emph{Recall:} say the Tamil for humiliation, then the Tamil for awkwardness, then say \emph{kūccam} again
\end{itemize}

\begin{tcolorbox}[breakable,colback=teal!4,colframe=teal!35!black,title={Before you move on}]
What does \ta{கூச்சம்} mean? (\textquotedblleft{}Bashfulness; ticklishness\textquotedblright{}.) Say it once more.
\end{tcolorbox}
\section[\texorpdfstring{\ta{அருவருப்பு} (aruvaruppu)}{அருவருப்பு (aruvaruppu)}]{\ta{அருவருப்பு} (aruvaruppu) — disgust}

\label{lesson:TA-C462-aruvaruppu}



\begin{quote}
Before the new one: say the Tamil for awkwardness, then the Tamil for bashfulness.
\end{quote}

\subsection*{You'll want to know: \ta{அருவருப்பு}}
\textbf{\ta{அருவருப்பு}} — \emph{aruvaruppu} — \textquotedblleft{}disgust\textquotedblright{}.

\begin{quote}
\ta{கெட்டுப்போன} \ta{உணவைப்} \ta{பார்த்து} \ta{எனக்கு} \ta{அருவருப்பாக} \ta{இருந்தது}.
\end{quote}

\begin{quote}
\emph{keṭṭuppōṉa uṇavaip pārttu eṉakku aruvaruppāka iruntatu.}
\end{quote}

\begin{quote}
The spoiled food disgusted me.
\end{quote}

\ta{அருவருப்பான}, disgusting.

\begin{cousinweb}
\textbf{\ta{அருவருப்பு}} (\emph{aruvaruppu}) is native Dravidian (DEDR 323), the noun of a verb \ta{அருவரு} \emph{aruvaru}, \textquotedblleft{}to loathe\textquotedblright{}.
\end{cousinweb}

\begin{grammarlens}[title={What you've built}]
One more word for describing what you lived and felt.
\end{grammarlens}

\subsection*{Your turn}
\begin{itemize}
\raggedright
  \item \emph{Say it:} \emph{aruvaruppu}
  \item \emph{Say it:} \emph{aruvaruppu}, once more
  \item \emph{Say it:} \emph{keṭṭuppōṉa uṇavaip pārttu eṉakku aruvaruppāka iruntatu.} — The spoiled food disgusted me.
  \item \emph{Recall:} say the Tamil for awkwardness, then the Tamil for bashfulness, then say \emph{aruvaruppu} again
\end{itemize}

\begin{tcolorbox}[breakable,colback=teal!4,colframe=teal!35!black,title={Before you move on}]
What does \ta{அருவருப்பு} mean? (\textquotedblleft{}Disgust\textquotedblright{}.) Say it once more.
\end{tcolorbox}
\section[\texorpdfstring{\ta{முகம்} \ta{சுளி} (mukam cuḷi)}{முகம் சுளி (mukam cuḷi)}]{\ta{முகம்} \ta{சுளி} (mukam cuḷi) — to grimace; to frown}

\label{lesson:TA-C462-mukam-culi}



\begin{quote}
Before the new one: say the Tamil for bashfulness, then the Tamil for disgust.
\end{quote}

\subsection*{You'll want to know: \ta{முகம்} \ta{சுளி}}
\textbf{\ta{முகம்} \ta{சுளி}} — \emph{mukam cuḷi} — \textquotedblleft{}to grimace; to frown\textquotedblright{}.

\begin{quote}
\ta{கசப்பான} \ta{மருந்தைக்} \ta{குடித்து} \ta{அவன்} \ta{முகம்} \ta{சுளித்தான்}.
\end{quote}

\begin{quote}
\emph{kacappāṉa maruntaik kuṭittu avaṉ mukam cuḷittāṉ.}
\end{quote}

\begin{quote}
He grimaced after drinking the bitter medicine.
\end{quote}

Past (I …): \ta{முகம்} \ta{சுளித்தேன்}, \emph{mukam cuḷittēṉ}.

\ta{முகம்} \textquotedblleft{}face\textquotedblright{} with \ta{சுளி} \textquotedblleft{}to show displeasure\textquotedblright{}.

\begin{cousinweb}
\textbf{\ta{முகம்} \ta{சுளி}} (\emph{mukam cuḷi}) is built from native words: \ta{முகம்} \emph{mukam}, \textquotedblleft{}face\textquotedblright{} (DEDR 4889), and \ta{சுளி} \emph{cuḷi}, \textquotedblleft{}to be angry, to show displeasure\textquotedblright{} (DEDR 2700).
\end{cousinweb}

\begin{grammarlens}[title={What you've built}]
One more word for describing what you lived and felt.
\end{grammarlens}

\subsection*{Your turn}
\begin{itemize}
\raggedright
  \item \emph{Say it:} \emph{mukam cuḷi}
  \item \emph{Say it:} \emph{mukam cuḷi}, once more
  \item \emph{Say it:} \emph{kacappāṉa maruntaik kuṭittu avaṉ mukam cuḷittāṉ.} — He grimaced after drinking the bitter medicine.
  \item \emph{Recall:} say the Tamil for bashfulness, then the Tamil for disgust, then say \emph{mukam cuḷi} again
\end{itemize}

\begin{tcolorbox}[breakable,colback=teal!4,colframe=teal!35!black,title={Before you move on}]
What does \ta{முகம்} \ta{சுளி} mean? (\textquotedblleft{}To grimace; to frown\textquotedblright{}.) Say it once more.
\end{tcolorbox}
